% Generated draft: requires mathematical and source review.
\documentclass[11pt]{article}
\usepackage[margin=25mm]{geometry}
\usepackage{fontspec}
\setmainfont{Latin Modern Roman}
\usepackage{amsmath,amssymb,mathtools,unicode-math}
\setmathfont{Latin Modern Math}
\usepackage{enumitem,needspace,xurl,hyperref}
\setlength{\parindent}{0pt}
\setlength{\parskip}{5pt}
\setlength{\emergencystretch}{3em}
\newcommand{\N}{\mathbb{N}}
\newcommand{\Z}{\mathbb{Z}}
\newcommand{\Q}{\mathbb{Q}}
\newcommand{\R}{\mathbb{R}}
\newcommand{\C}{\mathbb{C}}
\newcommand{\F}{\mathbb{F}}
\newcommand{\A}{\mathbb{A}}
\newcommand{\PP}{\mathbb{P}}
\newcommand{\E}{\mathbb{E}}
\newcommand{\eps}{\varepsilon}
\newcommand{\dd}{\,\mathrm{d}}
\begin{document}
\section*{UnsolvedMath ID: 30006603}
\textbf{Source problem code:} catalogue entry
\subsection{Definitions and mathematical statement}
Let \(K\) and \(L\) be simplicial complexes, with geometric realizations \(|K|\) and \(|L|\). They are \emph{PL-homeomorphic} if there is a homeomorphism between their realizations that is piecewise linear, with piecewise-linear inverse. A subdivision of a simplicial complex is a triangulation of its realization whose simplices lie within simplices of the original triangulation. A \emph{common refinement} of \(K\) and \(L\) is a simplicial complex \(H\) equipped with identifications under which it is a subdivision of each.

The extracted question asks whether every pair of PL-homeomorphic simplicial complexes has a common refinement obtainable from each by finitely many elementary refinement moves. To expose the unspecified hypotheses, let \(\mathcal C\) denote the intended class of complexes and let \(\mathcal M\) denote the intended family of elementary refinement moves. Write \(A\longrightarrow_{\mathcal M}B\) when one such move transforms \(A\) into \(B\). The proposed assertion has the form
\[
\begin{gathered}
\forall K,L\in\mathcal C,\qquad
|K|\cong_{\mathrm{PL}}|L|
\ \Longrightarrow\ 
\exists H\ \exists m,n\in\Z_{\geq 0}:\\
K=K_0\longrightarrow_{\mathcal M}K_1
\longrightarrow_{\mathcal M}\cdots
\longrightarrow_{\mathcal M}K_m\cong H,\\
L=L_0\longrightarrow_{\mathcal M}L_1
\longrightarrow_{\mathcal M}\cdots
\longrightarrow_{\mathcal M}L_n\cong H,
\end{gathered}
\]
where the terminal isomorphisms are simplicial and \(H\) is a common refinement. Such chains would express the corresponding relation as refinements followed by inverse refinements:
\[
K\longrightarrow_{\mathcal M}^{*}H
\longleftarrow_{\mathcal M}^{*}L.
\]
The source does not specify \(\mathcal C\), define the elementary moves, or explicitly state that the chains must be finite. Thus the displayed finite-chain formulation is a provisional interpretation requiring review. In particular, finiteness, local finiteness, dimension restrictions, and whether a prescribed PL homeomorphism must be realized are unresolved. The cited title concerning stellar subdivisions suggests a possible intended move family, but does not by itself establish the exact formulation intended here.

The original passage also asks the analogous question with merely homeomorphic realizations. This is a separate formulation, with \(|K|\cong|L|\) replacing \(|K|\cong_{\mathrm{PL}}|L|\). The supplied curation attributes the failure of the unrestricted homeomorphic-complex Hauptvermutung to Milnor, and attributes the modified PL version and the report's weighted strong-factorization result to Adiprasito--Pak. These are supplied literature assessments, not independently verified conclusions of this draft.

The passage further invokes Oda's strong factorization conjecture in toric geometry and states that a weighted version is true. It supplies no precise toric category, hypotheses, permitted operations, or definition of the weights. Consequently, no precise toric analogue or weighted factorization assertion can be recovered from this record alone.
\subsection{Short English statement}
The question asks whether PL-homeomorphic simplicial complexes admit a common subdivision reached from both by elementary refinement moves, yielding a factorization through that subdivision. The permitted complexes and moves are unspecified. The source also discusses the unrestricted homeomorphic version and toric analogues, and supplies a dated partially solved assessment; their exact scope requires human review.
\subsection{Sources}
\begin{itemize}
\item Local registry: \texttt{data/unsolvedmath/problems.json}, with \texttt{catalog/problems.json} as the supplementary registry; supplied complete record for permanent UnsolvedMath ID \(30006603\), source problem code \texttt{catalogue entry}, title \emph{Common Refinements of PL-Homeomorphic Simplicial Complexes}. Registry provenance is supplied; the local files were not inspected for this draft.
\item Supplied report citation: \emph{Median Geometry and Applications} (2026), Oberwolfach Reports, DOI \texttt{10.4171/owr/2026/8}, pp.\ 487--540. Literature-source label: Oberwolfach Report 8/2026. \url{https://doi.org/10.4171/owr/2026/8}. No exact page or speaker attribution for the extracted passage is supplied.
\item Supplied literature reference: Adiprasito--Pak, \emph{All triangulations have a common stellar subdivision}. \url{https://arxiv.org/abs/2404.05930}.
\item Literature assessment date: \texttt{2026-08-22}. This is supplied dated curation, not independently verified. The record labels the problem \texttt{partially\_solved}, attributes the disproof of the unrestricted homeomorphic-complex Hauptvermutung to Milnor, and attributes the modified PL version and the report's weighted strong-factorization result to Adiprasito--Pak. No specific Milnor reference or precise theorem statements are supplied. The background explicitly requests independent verification of time-sensitive status and cited claims.
\item Attribution and record qualifications: the collection is Oberwolfach Reports Open Problems; the proposed year is 2026 and no proposer is named. The supplied statement status is \texttt{corrected\_verified}, with the qualification that only formula delimiters, notation, grammar, and typography were repaired and mathematical content and scope were preserved. This metadata is not an independent mathematical verification.
\end{itemize}
\end{document}
